% =============================================================
%  数学建模竞赛 —— AI 工具使用说明 LaTeX 模板
%  编译方式: XeLaTeX（需 ctex 宏包支持中文）
%  使用说明: 将各处【】占位符替换为实际内容即可
% =============================================================
\documentclass[a4paper,12pt]{article}

% ---------- 中文支持 ----------
\usepackage[UTF8, scheme=plain]{ctex}

% ---------- 页面与版式 ----------
\usepackage{geometry}
\geometry{left=2.5cm, right=2.5cm, top=2.5cm, bottom=2.5cm}

\usepackage{setspace}
\onehalfspacing                              % 1.5 倍行距

% ---------- 字体与列表 ----------
\usepackage{enumitem}
\setlist{nosep, leftmargin=2em}              % 紧凑列表

% ---------- 超链接 ----------
\usepackage{hyperref}
\hypersetup{colorlinks=true, linkcolor=black, urlcolor=blue}

% ---------- 代码块 ----------
\usepackage{listings}
\usepackage{xcolor}
\lstset{
    basicstyle=\ttfamily\small,
    breaklines=true,
    frame=single,
    framesep=4pt,
    xleftmargin=1em,
    tabsize=4,
    showstringspaces=false,
}

% ---------- 自定义段落间距 ----------
\setlength{\parskip}{0.4em}

% =============================================================
\begin{document}

% ===================== 标题 =====================
\begin{center}
    {\zihao{3}\heiti AI 工具使用说明}
\end{center}
\vspace{0.5em}

\noindent\rule{\linewidth}{0.8pt}
\vspace{0.8em}

% ===================== 一、工具名称和版本 =====================
\section*{一、AI 工具名称和版本}
\noindent
【在此填写参赛使用的 AI 工具名称与版本，可列出多个工具，用顿号分隔。】

% 示例（可替换）：
% DeepSeek（V4 Pro）、通义千问（3.7 Max）

% ===================== 二、具体使用目的与环节 =====================
\section*{二、具体使用目的与环节}

% ---------- 工具 A ----------
\subsection*{\textbf{【AI 工具 A 名称（版本）】}}
\begin{enumerate}[label=(\arabic*)]
    \item \textbf{【使用环节一】}：【描述该环节中 AI 的具体使用目的与作用。】
    \item \textbf{【使用环节二】}：【描述第二个使用环节……】
    \item \textbf{【使用环节三】}：【描述第三个使用环节……】
\end{enumerate}

% ---------- 工具 B ----------
\subsection*{\textbf{【AI 工具 B 名称（版本）】}}
\begin{enumerate}[label=(\arabic*)]
    \item \textbf{【使用环节一】}：【描述该环节中 AI 的具体使用目的与作用。】
    \item \textbf{【使用环节二】}：【描述第二个使用环节……】
    \item \textbf{【使用环节三】}：【描述第三个使用环节……】
\end{enumerate}

% ===================== 三、关键交互记录 =====================
\section*{三、关键交互记录（重要提示词与回复）}

% ---------- 工具 A 交互记录 ----------
\subsection*{\textbf{【AI 工具 A 名称（版本）】}}

\noindent\textbf{Query 1：}\\
【填写第一条提示词原文。】

\noindent\textbf{Output：}\\
【填写 AI 的回复要点摘要。例如：分三大模块给出多组中英文检索关键词，区分图书、期刊文献检索方向……】

\vspace{0.6em}
\noindent\textbf{Query 2：}\\
【填写第二条提示词原文……】

\noindent\textbf{Output：}\\
【填写 AI 的回复要点摘要……】

\vspace{0.6em}
\noindent\textbf{Query 3：}\\
【填写第三条提示词原文……】

\noindent\textbf{Output：}\\
【填写 AI 的回复要点摘要……】

% ---------- 工具 B 交互记录 ----------
\subsection*{\textbf{【AI 工具 B 名称（版本）】}}

\noindent\textbf{Query 1：}\\
【填写第一条提示词原文……】

\noindent\textbf{Output：}\\
【填写 AI 的回复要点摘要……】

\vspace{0.6em}
\noindent\textbf{Query 2：}\\
【填写第二条提示词原文……】

\noindent\textbf{Output：}\\
【填写 AI 的回复要点摘要……】

\vspace{0.6em}
\noindent\textbf{Query 3：}\\
【填写第三条提示词原文……】

\noindent\textbf{Output：}\\
【填写 AI 的回复要点摘要……】

% ===================== 四、采纳和人工修改情况 =====================
\section*{四、采纳和人工修改情况}

% ---------- 工具 A ----------
\subsection*{\textbf{【AI 工具 A 名称（版本）】}}

\noindent\textbf{Query 1 采纳和人工修改情况说明：}\\
【说明采纳了 AI 哪些内容，人工做了哪些修改与核验。例如：采纳 AI 给出的检索关键词组合与综述基础写作框架；人工筛选贴合课题的内容，剔除无关领域解读，所有外文文献由团队手动进入数据库核对原始出版信息，修正 AI 给出的错误年份、页码。】

\vspace{0.6em}
\noindent\textbf{Query 2 采纳和人工修改情况说明：}\\
【说明采纳与人工修改情况……】

\vspace{0.6em}
\noindent\textbf{Query 3 采纳和人工修改情况说明：}\\
【说明采纳与人工修改情况……】

% ---------- 工具 B ----------
\subsection*{\textbf{【AI 工具 B 名称（版本）】}}

\noindent\textbf{Query 1 采纳和人工修改情况说明：}\\
【说明采纳与人工修改情况……】

\vspace{0.6em}
\noindent\textbf{Query 2 采纳和人工修改情况说明：}\\
【说明采纳与人工修改情况……】

\vspace{0.6em}
\noindent\textbf{Query 3 采纳和人工修改情况说明：}\\
【说明采纳与人工修改情况……】



\end{document}
